\chapter{Experimental Setup and Evaluation Metrics}
\label{chap:experiments}

\section{Dataset Specifications}

\subsection{Training Dataset}
\begin{itemize}
    \item Total scenarios: 50,000
    \item User counts: $N \in \{4, 6, 8, 10, 12\}$ (uniform distribution)
    \item Cell radius: 5000 m
    \item Carrier frequency: 3.5 GHz
    \item Total size: ~2.5 GB (CSV format)
\end{itemize}

\subsection{Data Split}
\begin{table}[H]
\centering
\begin{tabular}{lrr}
\toprule
\textbf{Split} & \textbf{Scenarios} & \textbf{Percentage} \\
\midrule
Training & 35,000 & 70\% \\
Validation & 7,500 & 15\% \\
Test & 7,500 & 15\% \\
\bottomrule
\end{tabular}
\caption{Dataset Split}
\end{table}

\subsection{Test Scenarios}

\textbf{In-Distribution:} Same as training ($N \in \{4,6,8,10,12\}$)

\textbf{Out-of-Distribution:}
\begin{enumerate}
    \item \textbf{Unseen Density:} $N=14$ users
    \item \textbf{Higher Shadowing:} $\sigma_{SF} = 8$ dB (vs 4-6 dB training)
    \item \textbf{Hexagonal Cells:} Different geometry
\end{enumerate}

\section{Hardware and Software Environment}

\begin{table}[H]
\centering
\begin{tabular}{ll}
\toprule
\textbf{Component} & \textbf{Specification} \\
\midrule
CPU & Intel Core i7-10700 @ 2.9 GHz \\
RAM & 16 GB DDR4 \\
GPU & NVIDIA GeForce RTX 3060 (6GB VRAM) \\
OS & Ubuntu 20.04 LTS \\
CUDA & 11.7 \\
Python & 3.9.13 \\
PyTorch & 1.13.0 \\
PyTorch Geometric & 2.2.0 \\
\bottomrule
\end{tabular}
\caption{Experimental Environment}
\end{table}

\section{Hyperparameters}

\subsection{GNN Architecture}
\begin{table}[H]
\centering
\begin{tabular}{lr}
\toprule
\textbf{Parameter} & \textbf{Value} \\
\midrule
Input features & 5 \\
Hidden dimension & 128 \\
Number of layers & 3 \\
Dropout rate & 0.3 \\
Aggregation & Mean \\
\bottomrule
\end{tabular}
\end{table}

\subsection{Training}
\begin{table}[H]
\centering
\begin{tabular}{lr}
\toprule
\textbf{Parameter} & \textbf{Value} \\
\midrule
Learning rate & $10^{-3}$ \\
Optimizer & Adam \\
Weight decay & $10^{-4}$ \\
Batch size & 32 \\
Max epochs & 200 \\
Early stopping patience & 25 \\
LR scheduler patience & 10 \\
LR reduction factor & 0.5 \\
Gradient clipping & 1.0 \\
\bottomrule
\end{tabular}
\end{table}

\subsection{Loss Weights}
\begin{table}[H]
\centering
\begin{tabular}{lr}
\toprule
\textbf{Task} & \textbf{Weight} \\
\midrule
Pairing (BCE) & 1.0 \\
Sum-Rate (MAE) & 0.5 \\
Power Allocation (MAE) & 0.3 \\
\bottomrule
\end{tabular}
\end{table}

\section{Evaluation Metrics}

\subsection{Classification Metrics (Pairing)}

\textbf{AUC-ROC:} Area under Receiver Operating Characteristic curve
\begin{equation}
\text{AUC} = \int_{0}^{1} \text{TPR}(\text{FPR}^{-1}(x)) dx
\end{equation}

\textbf{Precision:}
\begin{equation}
\text{Precision} = \frac{TP}{TP + FP}
\end{equation}

\textbf{Recall:}
\begin{equation}
\text{Recall} = \frac{TP}{TP + FN}
\end{equation}

\textbf{F1-Score:}
\begin{equation}
F1 = 2 \cdot \frac{\text{Precision} \times \text{Recall}}{\text{Precision} + \text{Recall}}
\end{equation}

\subsection{Regression Metrics (Sum-Rate, Power)}

\textbf{Mean Absolute Error (MAE):}
\begin{equation}
\text{MAE} = \frac{1}{n} \sum_{i=1}^{n} |\hat{y}_i - y_i|
\end{equation}

\textbf{Root Mean Squared Error (RMSE):}
\begin{equation}
\text{RMSE} = \sqrt{\frac{1}{n} \sum_{i=1}^{n} (\hat{y}_i - y_i)^2}
\end{equation}

\textbf{Mean Absolute Percentage Error (MAPE):}
\begin{equation}
\text{MAPE} = \frac{100\%}{n} \sum_{i=1}^{n} \left|\frac{\hat{y}_i - y_i}{y_i}\right|
\end{equation}

\textbf{R² Score:}
\begin{equation}
R^2 = 1 - \frac{\sum_{i}(\hat{y}_i - y_i)^2}{\sum_{i}(y_i - \bar{y})^2}
\end{equation}

\subsection{System Performance Metrics}

\textbf{System Throughput:}
\begin{equation}
R_{total} = \sum_{\text{all pairs}} R_{sum}(i,j)
\end{equation}

\textbf{Average User Rate:}
\begin{equation}
\bar{R} = \frac{R_{total}}{N_{paired}}
\end{equation}

\textbf{Throughput Ratio (vs Blossom):}
\begin{equation}
\eta = \frac{R_{total}^{\text{method}}}{R_{total}^{\text{Blossom}}} \times 100\%
\end{equation}

\textbf{Jain's Fairness Index:}
\begin{equation}
\mathcal{J} = \frac{\left(\sum_{i} R_i\right)^2}{N \sum_{i} R_i^2}
\end{equation}

\subsection{Computational Metrics}

\textbf{Runtime:} Wall-clock time for pairing (excluding power optimization)

\textbf{Speedup:}
\begin{equation}
\text{Speedup} = \frac{T_{\text{Blossom}}}{T_{\text{method}}}
\end{equation}

\textbf{Model Size:} Number of parameters and checkpoint file size

\section{Baseline Methods}

\begin{enumerate}
    \item \textbf{Static Pairing:} Strongest-with-weakest
    \item \textbf{Balanced Pairing:} Upper-half with lower-half
    \item \textbf{Blossom Matching:} Optimal general graph matching
    \item \textbf{Bipartite PF:} Bipartite with proportional fairness
    \item \textbf{NOMA-GNN:} Proposed method
\end{enumerate}

\section{Experimental Procedure}

\subsection{Training Phase}
\begin{enumerate}
    \item Generate 50k scenarios with 3GPP UMa channels
    \item Compute ground-truth labels via Blossom
    \item Split data 70/15/15
    \item Train NOMA-GNN for up to 200 epochs
    \item Save best model based on validation AUC
\end{enumerate}

\subsection{Testing Phase}
\begin{enumerate}
    \item Load test scenarios
    \item Run all 5 methods on each scenario
    \item Record throughput, runtime, fairness
    \item Compute aggregate statistics
\end{enumerate}

\subsection{Ablation Studies}
\begin{enumerate}
    \item Remove angular features: Train without $\theta_i$
    \item Vary encoder depth: 2, 3, 4 layers
    \item Single-task: Train only pairing classifier
    \item No hard negatives: Uniform negative sampling
    \item Different hidden dimensions: 64, 128, 256
\end{enumerate}

\subsection{Generalization Tests}
\begin{enumerate}
    \item Test on $N=14$ users (unseen density)
    \item Test with $\sigma_{SF} = 8$ dB shadowing
    \item Test on hexagonal cell geometry
\end{enumerate}

\section{Statistical Significance}

All results reported with:
\begin{itemize}
    \item Mean ± standard deviation over test set
    \item 95\% confidence intervals
    \item Paired t-tests for method comparisons ($p < 0.05$)
\end{itemize}

\section{Reproducibility Statement}

All experiments are reproducible via:
\begin{itemize}
    \item Fixed random seeds: 42
    \item Versioned dependencies: requirements.txt
    \item Saved configurations: config.py
    \item Public code repository: [to be added]
    \item Dataset generation scripts
\end{itemize}
